% !TEX root = Lecture_Koszul.tex
\clearpage
\section{Koszul complex and Koszul algebras}\label{sec:koszul-complex}

\subsection{The Koszul complex of a polynomial algebra}

Let $V$ be an $n$-dimensional $\Bbbk$-vector space, choose a basis $x_1,\ldots,x_n$, and let $A=\SymAlg(V)=\Bbbk[x_1,\ldots,x_n]$, with each $x_j$ in degree $1$.
Let $\pi:A\to\Bbbk$ be projection onto degree $0$, so $\Bbbk=A/A_+$.
\phantomsection\label{not:koszul-terms}
For $0\leq h\leq n$, put
\[
K_h:=\bigwedge^h V\otimes_\Bbbk A(-h),
\qquad K_0=A.
\]
The right $A$-action is on the second factor.
The factor $\bigwedge^hV$ is placed in grading degree $0$, so the generators of $K_h$ have grading degree $h$.

\begin{definition}\label{def:koszul-complex}
For $1\leq h\leq n$, define the \defn{Koszul differential}
\begin{equation}\label{eq:koszul-differential}
\begin{aligned}
d_h:K_h&\longrightarrow K_{h-1},\\
(v_1\wedge\cdots\wedge v_h)\otimes a
&\longmapsto
\sum_{r=1}^h(-1)^{r-1}
(v_1\wedge\cdots\wedge\widehat{v_r}\wedge\cdots\wedge v_h)
\otimes v_ra .
\end{aligned}
\end{equation}
The formula is multilinear and alternating in $v_1,\ldots,v_h$, so it defines a map on $\bigwedge^hV$.
The sequence $K_\bullet\to\Bbbk$ consists of these maps followed by $\pi:K_0=A\to\Bbbk$.
\end{definition}

\begin{proposition}\label{prop:koszul-complex-resolution}
The sequence
\[
0\longrightarrow K_n\xrightarrow{d_n}\cdots\xrightarrow{d_2}K_1
\xrightarrow{d_1}K_0\xrightarrow{\pi}\Bbbk\longrightarrow0
\]
is an exact minimal graded free resolution of $\Bbbk$ over $A$.
\end{proposition}

\begin{proof}
Each $d_h$ preserves degree and is right $A$-linear.
For $h\geq2$, compute $d_{h-1}d_h$ on a basis wedge.
Every term is indexed by two distinct positions $r<s$ deleted from that wedge.
Deleting position $r$ and then $s$ gives coefficient $(-1)^{r+s-3}x_{j_s}x_{j_r}$; deleting $s$ and then $r$ gives coefficient $(-1)^{r+s-2}x_{j_r}x_{j_s}$.
Since $x_{j_s}x_{j_r}=x_{j_r}x_{j_s}$, these two terms sum to zero. Hence $d_{h-1}d_h=0$; also $\pi d_1=0$ because every term of $d_1$ is divisible by a variable.

To prove exactness, give $A$ and $\bigwedge V$ the multigrading with $\deg x_j$ equal to the $j$th standard basis vector.
For $\alpha\in\mathbb Z_{\geq0}^n$, let $\operatorname{supp}(\alpha)=\{j\mid\alpha_j>0\}$.
The degree $\alpha$ part of $K_h$ has basis
\[
b_J:=(x_{j_1}\wedge\cdots\wedge x_{j_h})\otimes x^{\alpha-\mathbf1_J},
\qquad
J=\{j_1<\cdots<j_h\}\subseteq\operatorname{supp}(\alpha),
\]
where $\mathbf1_J$ has entry $1$ at indices in $J$ and $0$ elsewhere.
The differential is $d_h(b_J)=\sum_{r=1}^h(-1)^{r-1}b_{J\setminus\{j_r\}}$.

If $\alpha\neq0$, let $q$ be the least index in $\operatorname{supp}(\alpha)$.
For each $h\geq0$, define $s_{\alpha,h}:K_{h,\alpha}\to K_{h+1,\alpha}$ on basis vectors by
\[
s_{\alpha,h}(b_J)=
\begin{cases}
b_{\{q\}\cup J},&q\notin J,\\
0,&q\in J,
\end{cases}
\]
inserting $q$ first in the ordered wedge.
For $q\notin J$ and $h\geq1$, direct substitution gives
\[
\begin{aligned}
d_{h+1}s_{\alpha,h}(b_J)
&=b_J-\sum_{r=1}^h(-1)^{r-1}b_{\{q\}\cup(J\setminus\{j_r\})},\\
s_{\alpha,h-1}d_h(b_J)
&=\sum_{r=1}^h(-1)^{r-1}b_{\{q\}\cup(J\setminus\{j_r\})}.
\end{aligned}
\]
For $h=0$, $J=\varnothing$ and $d_1s_{\alpha,0}(b_\varnothing)=b_\varnothing$.
For $q\in J$, the term of $d_h(b_J)$ deleting $q$ is the only term on which $s_{\alpha,h-1}$ is nonzero, and it gives $s_{\alpha,h-1}d_h(b_J)=b_J$ while $d_{h+1}s_{\alpha,h}(b_J)=0$.
Thus $d_{h+1}s_{\alpha,h}+s_{\alpha,h-1}d_h=\id_{K_{h,\alpha}}$ for every $h\geq1$ in positive multidegree; the case $h=0$ was checked separately.
All positive multidegrees are therefore exact.
In multidegree $0$, only $K_0\cong\Bbbk$ is nonzero and $\pi$ is the identity onto $\Bbbk$.
This proves exactness of the displayed sequence.

Finally, each $K_h$ is a finite direct sum of copies of $A(-h)$, and \eqref{eq:koszul-differential} has image in $K_{h-1}A_+$.
Exactness gives $\Ker d_h=\Im d_{h+1}\subseteq K_hA_+$ for $h<n$, and $\Ker d_n=0$.
Also, $\Ker\pi=AA_+$.
Thus each map from $K_h$ onto the preceding kernel has kernel contained in $K_hA_+$, so it is a projective cover by Definition~\ref{def:graded-projective-cover}.
The resolution is therefore minimal.
\end{proof}

\subsection{Koszul algebras}

\begin{definition}\label{def:linear-resolution}
A minimal graded projective resolution $P_\bullet\to M$ is \defn{linear} if $P_h/P_hA_+$ is concentrated in degree $h$ for every $h\geq0$.
This is equivalent to $P_h$ being generated in degree $h$: generation in degree $h$ implies the stated concentration because $A_+$ raises degree.
Conversely, if the top is concentrated in degree $h$, boundedness below and Lemma~\ref{lem:graded-nakayama} show that $P_h$ has no components below degree $h$; the quotient $P_h/(P_h)_hA$ then has zero graded top, so Lemma~\ref{lem:graded-nakayama} makes that quotient zero.
\end{definition}

\begin{definition}\label{def:koszul-algebra}
Let $A=T_S(V)/(R)$ be a quadratic algebra with $S=A_0$ finite-dimensional semisimple.
The algebra $A$ is \defn{Koszul} if the degree zero quotient module $S=A/A_+$ has a linear minimal graded projective resolution as a right $A$-module.
\end{definition}

The generators in homological degree $h$ must then all have grading degree $h$.
Since $S$ can have several simple modules, the definition requires a linear resolution of their direct sum, not just of one chosen simple module.

\begin{proposition}\label{prop:polynomial-koszul}
For every finite-dimensional vector space $V$, the symmetric algebra $\SymAlg(V)$ is Koszul.
\end{proposition}

\begin{proof}
Proposition~\ref{prop:koszul-complex-resolution} gives a minimal resolution with $K_h=\bigwedge^hV\otimes_\Bbbk A(-h)$.
Thus $K_h$ is generated in degree $h$ for every $h$.
\end{proof}

\begin{example}\label{ex:koszul-nonconnected}
The definition also applies when $S$ has several simple modules.
For $A=\Bbbk A_2$, the resolution in Example~\ref{ex:minimal-A2} is linear for $L_1$, and $L_2=e_{22}A$ is projective in degree $0$.
Taking their direct sum gives a linear resolution of $S=L_1\oplus L_2$; hence $\Bbbk A_2$ is Koszul.
\end{example}

\begin{example}
The exterior algebra $E=\Bbbk[\epsilon]/(\epsilon^2)$ is quadratic, and Example~\ref{ex:minimal-exterior} gives a minimal resolution of its degree zero quotient with $P_h=E(-h)$.
Thus $E$ is Koszul, although its minimal resolution is infinite.
\end{example}


\subsection{Exercises}

\begin{exercise}
Let $A=\SymAlg(V)=\Bbbk[x_1,\ldots,x_n]$ and $K_\bullet\to\Bbbk$ be the Koszul resolution of Proposition~\ref{prop:koszul-complex-resolution}.
In each degree, the alternating sum of the dimensions in an exact sequence of finite-dimensional vector spaces is zero.
Use this to show
\[
(1-t)^n\,h_A(t)=1,
\]
and deduce the Hilbert series of $\Bbbk[x_1,\ldots,x_n]$.
\end{exercise}

\begin{exercise}
Let $S=\bigoplus_{y=1}^n\Bbbk e_y$, let $V$ be a finite-dimensional bimodule in degree $1$, and let $A=T_S(V)/(V\otimes_SV)=S\oplus V$.
Prove that the minimal graded projective resolution of $S$ has $P_h\cong(V^{\otimes_Sh}\otimes_SA)(-h)$ for all $h\geq0$, and deduce that $A$ is Koszul.

\emph{Hint.} Show that $\Ker p_0=V$ and that $VA_+=0$, so each kernel is killed by $A_+$.
\end{exercise}

\begin{exercise}
Let $A=\Bbbk[x]/(x^3)$ with $\deg x=1$.
Find the minimal graded projective resolution of $\Bbbk$, and show that $P_h$ is generated in degree $\lfloor 3h/2\rfloor$ for all $h\geq0$.
Deduce that $A$ is not Koszul.
Is $A$ quadratic?
\end{exercise}
